
\bigskip

\section{Hyperimaginaries}\label{hyperimaginaries}
In this section, we review some basics about hyperimaginaries. As we said in the Introduction, hyperimaginaries are important for the study of simple unstable theories. For standard textbook references, see \cites[\S3\addsemicolon]{alma992983361725701631}[\S\S15, 18]{casanovas2011simple}. 

\begin{defin}
    A theory $T$ \textit{eliminates hyperimaginaries} if for every type-definable equivalence relation $E$ on a type $X$, there is a sequence of definable equivalence relations $E_i$ on definable sets $X_i$, such that $\bigcap_i E_i \upharpoonright_X = E$.
\end{defin}

The following two propositions show that instead of working with $T$, we may work over $T^{eq}$.

\begin{prop}\label{EHI invariant under eq}
    A theory $T$ eliminates hyperimaginaries if and only if $T^{eq}$ also eliminates hyperimaginaries.
\end{prop}

\begin{proof}
    $(\Rightarrow)$ \quad  Let $E((y_i)_{i\in I},(y'_i)_{i\in I})$ be a type-definable equivalence relation in $T^{eq}$, where $y_i$ is of sort $\M/R_i$. To simplify notation, we write $y_I$ for $(y_i)_{i\in I}$. Then, $E(\pi_I(x_I),\pi_I(x'_I))$ is a type-definable equivalence relation in $T$, where $\pi_I$ is the product of the quotient maps $\pi_i: \M\rightarrow \M/R_i$. Since $T$ eliminates hyperimaginaries, there are definable equivalence relations $E_J(x_J,x'_J)$ in $T$, where $J\in\mathcal{J}$ is a finite subset of $I$, such that for all $a_I,a'_I\in\M$,
    \begin{equation}
        E(\pi_I(a_I),\pi_I(a'_I)) \ \Leftrightarrow \ E_J(a_J,a'_J) \textup{ for all }J\in\mathcal{J}. \tag{$*$}
    \end{equation}  
    For each $J$, consider the definable relation
    \[ \varepsilon_J(y_J,y'_J) := \exists x_J, x'_J \ (\pi_J(x_J)=y_J \wedge \pi_J(x'_J)=y'_J \wedge E_J(x_J, x'_J)). \]
    The reflexivity and symmetry of $\varepsilon_J$ is easy to see. 
    For transitivity, we need the following.

    \begin{claim}\label{EJ RJ}
        Suppose $\pi_J(a_J)=\pi_J(a'_J)$ for some $J\in\mathcal{J}$. Then, $E_J(a_J,a'_J)$ holds.
    \end{claim}

    \begin{proof}
        Suppose $\pi_J(a_J)=\pi_J(a'_J)$ holds. Extend $a_J$ in any way to a tuple $a_I$ and extend $a'_J$ to $a'_I$ such that $a_i = a'_i$ for all $i\not\in J$. Then, we have $\pi_I(a_I)=\pi_I(a'_I)$, so that by the reflexivity of $E$, we have $E(\pi_I(a_I),\pi_I(a'_I))$. By $(\ast)$, we have $E_J(a_J,a'_J)$. 
    \end{proof}
    
    Back to the proof of the transitivity of $\varepsilon_J$. Suppose we have $\varepsilon_J(b_J,b'_J)$ and $\varepsilon_J(b'_J,b''_J)$. By definition, there are $a_J,a'_J,\alpha'_J,$ and $\alpha''_J$ such that 
    \begin{align*}
        \pi_J(a_J) &=b_J &\wedge& &\pi_J(a'_J) &=b'_J &\wedge & &E_J(a_J,a'_J); \\
        \pi_J(\alpha'_J) &= b'_J & \wedge & &\pi_J(\alpha''_J) &=b''_J &\wedge &&E_J(\alpha'_J,\alpha''_J).
    \end{align*}
    By the Claim, we have $E_J(a'_J,\alpha'_J)$, so the transitivity of $E_J$ implies that $E_J(a_J,\alpha''_J)$. Then, $a_J,\alpha''_J$ witnesses $\varepsilon_J(b_J,b''_J)$ as required.
    
    Finally, we prove that $\bigcap_{J\in\mathcal{J}} \varepsilon_J = E$. Fix any $b_i,b'_i$ of sort $\M/R_i$, and any $a_i,a'_i$ such that $\pi_i(a_i)=b_i$ and $\pi_i(a'_i)=b'_i$ for all $i\in I$. Note that if we have $\varepsilon_J(b_J,b'_J)$, so that there are some $\alpha_J,\alpha'_J$ such that $\pi_J(\alpha_J)=b_J$ and $\pi_J(\alpha'_J)=b'_J$ satisfying $E_J(\alpha_J,\alpha'_J)$, by the Claim, we also have $E_J(a_J,a'_J)$. Hence, we conclude that
    \[\varepsilon_J(b_J,b'_J) \textup{ for all } J\in\mathcal{J} \ \Leftrightarrow \ E_J(a_J,a'_J) \textup{ for all } J\in\mathcal{J} \ \Leftrightarrow \ E(b_I,b'_I), \]
    where the second equivalence comes from $(*)$.

\smallskip

    $(\Leftarrow)$ \quad  Suppose $E$ is a type-definable equivalence relation in $T$. Then, there are definable equivalence relations $E_i$ in $T^{eq}$ such that $\bigcap E_i = E$. Now, $E_i$ is a definable relation on the home sort, so in fact each $E_i$ is definable in $T$. 
\end{proof}

\begin{prop}\label{if prodef exact then def exact}
    If $\Pro(\defT(T))$ is exact, then $\defT(T)$ is exact.
\end{prop}

\begin{proof}
    Suppose $\Pro(\defT(T))$ is exact. Fix an equivalence relation $E\rightrightarrows X$ in $\defT(T)$. Then, there is a quotient $f: X\rightarrow Y$ in $\Pro(\defT(T))$, represented by some $f_i: X\rightarrow Y_i$ in $\defT(T)$. Logically, $f$ being the quotient map for $E$ means 
    \[ E(x,x') \ \Leftrightarrow \ f_i(x)=f_i(x') \textup{ for all }i. \]
    By compactness, there exists an $i_0$ such that $E(x,x')$ is equivalent to $f_{i_0}(x)=f_{i_0}(x')$. Then, $f_{i_0}:X\rightarrow Y_{i_0}$ is the quotient for $E$ in $\defT(T)$.
\end{proof}

\begin{thm}\label{eliminate hyper iff pro exact}
    A theory $T$ eliminates hyperimaginaries if and only if $\Pro(\defT(T^{eq}))$ is exact.
\end{thm}

\begin{proof}
    By Proposition \ref{EHI invariant under eq}, we assume $T$ eliminates imaginaries.

\smallskip

    $(\Rightarrow)$ \quad Let $E\rightrightarrows X$ be an equivalence relation in $\Pro(\defT(T))$. Let $(E_i\rightrightarrows X_i)_{i\in I}$ be a sequence of definable equivalence relations which eliminate $E$. Without loss of generality, assume $I$ is a cofiltered system. As $T$ eliminates imaginaries, the quotients $X_i/E_i$ are definable. Given a map $i\rightarrow j$ in $I$,  the universality of $X_i/E_i$ as a coequaliser induces a map $X_i/E_i\rightarrow X_j/E_j$. 
% https://q.uiver.app/#q=WzAsNixbMCwwLCJFX2kiXSxbMSwwLCJYX2kiXSxbMCwxLCJFX2oiXSxbMSwxLCJYX2oiXSxbMiwwLCJYX2kvRV9pIl0sWzIsMSwiWF9qL0VfaiJdLFswLDEsIiIsMSx7Im9mZnNldCI6MX1dLFswLDEsIiIsMSx7Im9mZnNldCI6LTF9XSxbMiwzLCIiLDEseyJvZmZzZXQiOjF9XSxbMiwzLCIiLDEseyJvZmZzZXQiOi0xfV0sWzAsMl0sWzEsM10sWzEsNCwiZV9pIl0sWzMsNSwiZV9qIiwyXSxbNCw1LCIiLDEseyJzdHlsZSI6eyJib2R5Ijp7Im5hbWUiOiJkYXNoZWQifX19XV0=
\[\begin{tikzcd}
	{E_i} & {X_i} & {X_i/E_i} \\
	{E_j} & {X_j} & {X_j/E_j}
	\arrow[shift right, from=1-1, to=1-2]
	\arrow[shift left, from=1-1, to=1-2]
	\arrow[from=1-1, to=2-1]
	\arrow["{e_i}", from=1-2, to=1-3]
	\arrow[from=1-2, to=2-2]
	\arrow[dashed, from=1-3, to=2-3]
	\arrow[shift right, from=2-1, to=2-2]
	\arrow[shift left, from=2-1, to=2-2]
	\arrow["{e_j}"', from=2-2, to=2-3]
\end{tikzcd}\]

    Hence, we obtain the cofiltered limit $\lim X_i/E_i \in \Pro(\defT(T))$. Moreover, the maps $e_i:X_i \rightarrow X_i/E_i$ induce a map $e=\lim e_i:X\rightarrow \lim X_i/E_i$. By Lemma \ref{level regular epimorphism}, $e$ is a regular epimorphism. Kernel pairs commute with cofiltered limits, so we have
    \[ X\times_{\lim X_i/E_i} X 
    \cong \lim (X_i\times_{X_i/E_i}X_i) 
    \cong \lim E_i
    \cong E.\]
    Hence, $E\rightrightarrows X$ is a kernel pair with coequaliser $\lim X_i/E_i$ as required.

    \smallskip

    $(\Leftarrow)$ \quad Let $E\rightrightarrows X$ be a type-definable equivalence relation. By exactness, there is a type-definable quotient $Y = X/E$. By Lemma \ref{level intersections}, the map $\pi:X\rightarrow X/E$ is the intersection of some $(\pi_i:X_i\rightarrow Y_i)$, where $X=\bigcap X_i$ and $Y=\bigcap Y_i$.

    Let $E_i$ be the pullback $X_i \times_{Y_i} X_i$. It is clear that $E_i$ is a definable equivalence relation. Moreover, in a saturated model $\M$ for every $a,b\in \M(X)$, we have:
    \begin{align*}
        E(a,b) &\Leftrightarrow \pi(a)=\pi(b) \\
        &\Leftrightarrow \pi_i(a)=\pi_i(b)  \textup{ for all } i \\
        &\Leftrightarrow E_i(a,b) \textup{ for all } i.
    \end{align*}
    Hence, $\M(E)=\bigcap \M(E_i)$. Since $T$ is complete, $E=\bigcap E_i$.
\end{proof}

Noting that the proof of Theorem \ref{eliminate hyper iff pro exact} does not use the fact that $\defT(T)$ is boolean, we in fact obtain a more general theorem: the pro-completion of an exact category $\C$ is exact if and only if every equivalence relation in $\Pro(\C)$ is the cofiltered limit of equivalence relations in $\C$. The latter condition was studied in \cite[Ch.~III \S4]{meyerphd}. Specifically, it is proved that reflexive symmetric relations in $\Pro(\C)$ are cofiltered limits of reflexive symmetric relations in $\C$. This is in fact a corollary of uniform approximation. In the notation of Theorem \ref{simonhenry}, let $I$ be the following category
    % https://q.uiver.app/#q=WzAsMixbMCwwLCJSIl0sWzEsMCwiWCwiXSxbMSwwLCJyIiwxXSxbMCwxLCJwXzEiLDEseyJvZmZzZXQiOi0zfV0sWzAsMSwicF8yIiwxLHsib2Zmc2V0IjozfV0sWzAsMCwicyIsMV1d
    \[\begin{tikzcd}
    	R & {X,}
    	\arrow["s"{description}, from=1-1, to=1-1, loop, in=55, out=125, distance=10mm]
    	\arrow["{p_1}"{description}, shift left=3, from=1-1, to=1-2]
    	\arrow["{p_2}"{description}, shift right=3, from=1-1, to=1-2]
    	\arrow["r"{description}, from=1-2, to=1-1]
    \end{tikzcd}\]
where $r$ satisfies the usual property of reflexivity and $s$ satisfies the usual property of symmetry. Now, $I$ has a non-identity endomorphism $s$, but part (1) of Theorem \ref{simonhenry} still applies as $\defT(T)$ has finite limits.

However, one cannot do this for transitivity. Indeed, the diagram for an equivalence relation is the following:
    % https://q.uiver.app/#q=WzAsMyxbMSwwLCJSIl0sWzIsMCwiWCwiXSxbMCwwLCJSXFx0aW1lc19YIFIiXSxbMSwwLCJyIiwxXSxbMCwxLCJwXzEiLDEseyJvZmZzZXQiOi0zfV0sWzAsMSwicF8yIiwxLHsib2Zmc2V0IjozfV0sWzAsMCwicyIsMV0sWzIsMCwidCIsMV1d
    \[\begin{tikzcd}
    	{R\times_X R} & R & {X,}
    	\arrow["t"{description}, from=1-1, to=1-2]
    	\arrow["s"{description}, from=1-2, to=1-2, loop, in=55, out=125, distance=10mm]
    	\arrow["{p_1}"{description}, shift left=3, from=1-2, to=1-3]
    	\arrow["{p_2}"{description}, shift right=3, from=1-2, to=1-3]
    	\arrow["r"{description}, from=1-3, to=1-2]
    \end{tikzcd}\]
where $r$ and $s$ are as before, and $t$ satisfies the usual property of transitivity. Of course, one can let $I$ to be the category above and apply Theorem \ref{simonhenry}. Then, one obtains morphisms $t_i:W_i\rightarrow E_i$, whose cofiltered limit is $t$, but there is no guarantee that $W_i \cong E_i\times_{X_i} E_i$, or that the relations $E_i\rightrightarrows X_i$ are equivalence relations.

\smallskip

The proof of Theorem \ref{eliminate hyper iff pro exact} provides the precise construction of quotients in $\Pro(\defT(T))$. We apply this to the example of the infinitary Chinese remainder theorem with $T=\textup{Th}(\mathbb{Z},+,\times,0,1)$. We are interested in the quotient for the type-definable equivalence relation $\bigcap_{p \textup{ prime}} \textup{mod }p$. While $T$ might not eliminate all hyperimaginaries, $T$ certainly eliminates this particular one, which is already given as the intersection of definable equivalence relations. Hence, the quotient is also type-definable.

The quotient is constructed as the cofiltered limit of the quotients of the equivalence relations $\bigcap_{p\in I} \textup{mod } p$, where $I$ is a finite set of prime numbers.
The ordinary Chinese remainder theorem states that finite products and quotients commute: $\prod_{p\in I}\mathbb{Z}/p\mathbb{Z} \cong \mathbb{Z}/(\bigcap_{p\in I}\textup{mod }p)$. Hence, the quotient of $\bigcap_{p \textup{ prime}} \textup{mod }p$ is the profinite integers $\prod_{p \textup{ prime}}\mathbb{Z}/p\mathbb{Z}$.

This might seem surprising. Consider that $\bigcap_{p \textup{ prime}} \textup{mod }p$ is the trivial equivalence relation on $\mathbb{Z}$, with quotient $\mathbb{Z}$ itself, not $\prod_{p \textup{ prime}}\mathbb{Z}/p\mathbb{Z}$. The problem is that one should not calculate quotients inside the standard model $\mathbb{Z}$. Indeed, given a sufficiently saturated model $\M$, we do have $\M/(\bigcap_{p \textup{ prime}} \textup{mod }p) \cong \prod_{p \textup{ prime}} \mathbb{Z}/p\mathbb{Z}$. 

\begin{example}
    It was noted in \cite[\S1]{Pillay_Poizat_1987} that the theory \textsf{RCF} of real closed fields does not eliminate hyperimaginaries. In particular, the type-definable equivalence relation 
    \[ E(x,y):= \{|x-y| < \frac{1}{n} \ | \ n\in\mathbb{Z}^+ \}\]
    of being infinitesimally close together is not eliminated. The original argument uses o-minimality, but Theorem \ref{eliminate hyper iff pro exact} yields a new argument. Noting that \textsf{RCF} eliminates imaginaries, suppose \textsf{RCF} also eliminates hyperimaginaries, so that $\Pro(\defT(\textsf{RCF}))$ is exact. Then, the quotient of $E$ is type-definable, say by a type $X$. Now, an $E$-hyperimaginary is a set of numbers infinitesimally close to each other and is represented by their shared standard part. Hence, $X$ interpreted by a saturated model $\M$ is isomorphic to the standard $\mathbb{R}$ and does not contain an infinitesimal. However, one can force infinitesimals to exist in a saturated model by realising the type $X\cup\{0<x<\frac{1}{n}\ | \ n\in\mathbb{Z}^+ \}$, yielding a contradiction.
\end{example}

We draw a corollary from Theorem~\ref{eliminate hyper iff pro exact}.

\begin{cor}
    A theory $T$ eliminates imaginaries and hyperimaginaries if and only if $\Pro(\defT(T))$ is exact.
\end{cor}

\begin{proof}
    $(\Rightarrow)$ \quad By Proposition~\ref{EI iff exact} and Theorems~\ref{Teq is exact} and ~\ref{eliminate hyper iff pro exact}. 
    
    \smallskip
    
    $(\Leftarrow)$ \quad By Proposition~\ref{if prodef exact then def exact}, we know $\defT(T)$ is exact, so by Proposition~\ref{EI iff exact}, $T$ eliminates imaginaries. By Theorem~\ref{Teq is exact}, we know $\defT(T) \simeq \defT(T)_{ex/reg} \simeq \defT(T^{eq})$. Hence, we have $\Pro(\defT(T^{eq})) \simeq \Pro(\defT(T))$, which is exact, so by Theorem~\ref{eliminate hyper iff pro exact}, $T$ eliminates hyperimaginaries.
    
\end{proof}
